\subsection{Fusion Categories}
\label{subsec:fusion-categories}

\begin{definition}[Premodular fusion category]
\label{def:premodular-category}
\begin{enumerate}[label=(\roman*),ref=\roman*]
  \item\label{item:fusion-category}
  A \emph{fusion category} is a finite semisimple $\mathbb C$-linear abelian
  rigid monoidal category whose tensor product is bilinear and whose tensor
  unit $\mathbf 1$ is simple.

  \item\label{item:ribbon-structure}
  A \emph{ribbon structure} on a braided fusion category is a natural
  automorphism $\theta$ of the identity satisfying
  \[
    \theta_{X\otimes Y}
    =(\theta_X\otimes\theta_Y)c_{Y,X}c_{X,Y},
    \qquad
    \theta_{X^*}=(\theta_X)^*.
  \]

  \item\label{item:premodular-category}
  A \emph{premodular fusion category} is a ribbon fusion category.
\end{enumerate}
\end{definition}

\begin{remark}[Terminology and references]
\label{rem:premodular-category-reference}
These are \cite[Definitions~4.1.1, 8.10.1, and 8.13.1]{EGNO2015}.
In particular, premodularity imposes no nondegeneracy condition on the
braiding.
\end{remark}

\begin{convention}[Simple-object data]
\label{conv:simple-object-data}
Fix a premodular fusion category $\cC$.
\begin{enumerate}[label=(\roman*),ref=\roman*]
  \item\label{item:simple-skeleton}
  The finite set $\operatorname{Irr}(\cC)$ contains one representative of
  each simple isomorphism class.

  \item\label{item:quantum-dimension}
  For $i\in\operatorname{Irr}(\cC)$, its spherical quantum dimension is
  denoted $d_i$, and its twist scalar is
  $\theta_i\in\mathbb C^\times$.

  \item\label{item:global-dimension}
  The global dimension is
  \[
    \operatorname{Dim}(\cC)
    :=\sum_{i\in\operatorname{Irr}(\cC)}d_i^2.
  \]
\end{enumerate}
\end{convention}

\begin{remark}[Dimension convention]
\label{rem:global-dimension-reference}
The notation and global-dimension convention are those of
\cite[Sections~4.1 and~8.13, especially Equation~(8.53)]{EGNO2015}.
\end{remark}

\begin{proposition}[Canonical spherical structure and dimension]
\label{prop:canonical-spherical-dimension}
Let \(\cC\) be a premodular fusion category.  Its ribbon structure determines
its canonical spherical pivotal
structure.  With respect to this structure, the global dimension in
\cref{item:global-dimension} is nonzero.
\end{proposition}

\begin{proof}
The first assertion is \cite[Proposition~8.10.12]{EGNO2015}.  The displayed
global-dimension formula is \cite[Equation~(8.53)]{EGNO2015}, and its
nonvanishing in characteristic zero follows from
\cite[Theorem~7.21.12]{EGNO2015}.
\end{proof}

\begin{definition}[Character]
\label{def:character}
Let \(\cC\) be a premodular fusion category.  A \emph{character} of \(\cC\)
is a monoidal natural automorphism
\[
  \eta\colon\operatorname{Id}_{\cC}\Longrightarrow
  \operatorname{Id}_{\cC}.
\]
\end{definition}

\begin{lemma}[Simple-object character scalars]
\label{lem:simple-object-character-scalars}
Let \(\eta\) be a character of a premodular fusion category \(\cC\).  For
every \(i\in\operatorname{Irr}(\cC)\), there is a unique
\(\eta_i\in\mathbb C^\times\) such that
\[
  (\eta)_i=\eta_i\operatorname{id}_i.
\]
\end{lemma}

\begin{proof}
Schur's lemma identifies \(\operatorname{End}_{\cC}(i)\) with
\(\mathbb C\).  Since \((\eta)_i\) is an automorphism, its scalar is
nonzero.
\end{proof}

\begin{remark}[Yetter's terminology]
\label{rem:yetter-character-terminology}
This is the ``functorial tensor automorphism'' of
\cite[Definition~1]{Yetter1993}.
\end{remark}

\begin{lemma}[Character identities]
\label{lem:character-identities}
Let \(\cC\) be a premodular fusion category and let \(\eta\) be a character
of \(\cC\).  For $i,j,k\in\operatorname{Irr}(\cC)$,
\begin{enumerate}[label=(\roman*),ref=\roman*]
  \item\label{item:character-unit-dual}
  $\eta_{\mathbf 1}=1$ and $\eta_{i^*}=\eta_i^{-1}$.

  \item\label{item:character-fusion}
  If $k$ is a simple summand of $i\otimes j$, then
  $\eta_k=\eta_i\eta_j$.
\end{enumerate}
\end{lemma}

\begin{proof}
The unit and tensor identities are the monoidal axioms.  Naturality for the
evaluation $i^*\otimes i\to\mathbf 1$ gives
$\eta_{i^*}\eta_i=1$.  If $u\colon k\to i\otimes j$ is a nonzero inclusion,
naturality gives $(\eta_i\otimes\eta_j)u=u\eta_k$, and Schur's lemma gives
the second identity.
\end{proof}

\begin{definition}[Universal grading group]
\label{def:universal-grading-group}
Let \(\cC\) be a fusion category.  The \emph{universal grading group}
\(U(\cC)\) is a group equipped with a map
\[
  \deg\colon\operatorname{Irr}(\cC)\longrightarrow U(\cC)
\]
such that
\[
  k\text{ is a summand of }i\otimes j
  \quad\Longrightarrow\quad
  \deg(k)=\deg(i)\deg(j),
\]
for all \(i,j,k\in\operatorname{Irr}(\cC)\), and which is universal among
groups equipped with such a map.
\end{definition}

\begin{remark}[Yetter's grading group]
\label{rem:yetter-grading-group}
This is the group denoted $\operatorname{Gr}_{\cC}$ by
\cite[Definition~3]{Yetter1993}.
\end{remark}

\begin{proposition}[Characters and the universal grading]
\label{prop:characters-universal-grading}
Let \(\cC\) be a fusion category.  There is a canonical isomorphism of
abelian groups
\[
  A_{\cC}:=\operatorname{Aut}_{\otimes}
    (\operatorname{Id}_{\cC})
  \cong
  \operatorname{Hom}\bigl(U(\cC),\mathbb C^\times\bigr).
\]
Under this isomorphism, $a\in A_{\cC}$ acts on
\(i\in\operatorname{Irr}(\cC)\) by
the scalar obtained by evaluating the corresponding character at
$\deg(i)$.
\end{proposition}

\begin{proof}
This is \cite[Proposition~4]{Yetter1993}; in modern terminology it is
\cite[Proposition~4.14.3]{EGNO2015}.
\end{proof}

\begin{remark}[Full and cyclic twists]
\label{rem:full-cyclic-twists}
The full Yetter coefficient group is $A_{\cC}$.  Fixing one
$\eta\in A_{\cC}$ gives a homomorphism
\[
  \mathbb Z\longrightarrow A_{\cC},\qquad n\longmapsto\eta^n,
\]
and hence a cyclic specialization of the full theory.  This paper proves the
fixed-character specialization with $\alpha\in H_2(-;\mathbb Z)$.  The full
$H_2(-;A_{\cC})$-valued extension is discussed in
\cref{sec:future-directions}; it is not claimed as a theorem here.
\end{remark}

\begin{definition}[$\cC$-colored ribbon graph]
\label{def:colored-ribbon-graph}
Let \(\cC\) be a premodular fusion category.  A $\cC$-colored ribbon graph
in an oriented $3$-ball consists of the
following data, considered up to ambient isotopy relative to the boundary.
\begin{enumerate}[label=(\roman*),ref=\roman*]
  \item\label{item:ribbon-embedding}
  A compact oriented framed graph properly embedded in the ball.

  \item\label{item:ribbon-labels}
  An object of $\cC$ labelling every oriented ribbon.  Reversing a ribbon
  orientation replaces its label $X$ by $X^*$.

  \item\label{item:ribbon-coupons}
  An actual morphism labelling every coupon.  A coupon with ordered incoming
  labels $X_1,\ldots,X_r$ and ordered outgoing labels
  $Y_1,\ldots,Y_s$ carries a morphism
  \[
    f\colon X_1\otimes\cdots\otimes X_r
    \longrightarrow Y_1\otimes\cdots\otimes Y_s.
  \]

  \item\label{item:ribbon-boundary}
  Designated input and output disks containing the ordered incoming and
  outgoing boundary ribbons.
\end{enumerate}
\end{definition}

\begin{convention}[Graphical generators]
\label{conv:graphical-generators}
Let \(\cC\) be a premodular fusion category and let \(\eta\) be a character
of \(\cC\).  In a generic planar projection,
crossings, cups, caps, and full twists denote the braiding, duality maps, and
ribbon twist.  A character coupon on an $X$-colored ribbon denotes
$\eta_X^n$ for a specified $n\in\mathbb Z$.
\end{convention}

\begin{remark}[Ribbon-graph references]
\label{rem:ribbon-graph-references}
This is the standard coloured-ribbon-graph language of
\cite[Chapter~2, Sections~2.2--2.3]{BakalovKirillov2001}; the corresponding
free tortile category is treated in \cite[Theorem~6.1]{Shum1994}.
\end{remark}

\begin{proposition}[Graphical interpretation]
\label{prop:graphical-interpretation}
Let \(\cC\) be a premodular fusion category.  A $\cC$-colored ribbon graph
in a $3$-ball canonically determines a morphism
from the tensor product of its ordered input labels to the tensor product of
its ordered output labels.  This morphism is invariant under ambient isotopy
relative to the boundary.
\end{proposition}

\begin{proof}
Choose a generic planar projection, interpret horizontal juxtaposition as
tensor product and vertical stacking as composition, and apply the ribbon
relations.  Independence of all choices is the graphical calculus of
\cite[Chapter~2, especially Sections~2.2--2.3]{BakalovKirillov2001}; see also
the ribbon coherence theorem of \cite{Shum1994}.
\end{proof}

\begin{definition}[Reshetikhin--Turaev graph evaluation]
\label{def:rt-graph-evaluation}
Let \(\cC\) be a premodular fusion category and let
$\Gamma\subset S^3$ be a closed $\cC$-colored ribbon graph.  Remove a
$3$-ball disjoint from $\Gamma$ and apply
\cref{prop:graphical-interpretation}.  The resulting endomorphism of
$\mathbf 1$ is a scalar.  We denote it by
\[
  \operatorname{RT}_{\cC}(S^3,\Gamma)\in\mathbb C.
\]
\end{definition}

\begin{remark}[Closed-graph specialization]
\label{rem:rt-closed-graph-reference}
This is the closed-graph specialization of the Reshetikhin--Turaev graphical
functor in \cite[Chapter~2, Section~2.3]{BakalovKirillov2001}.
\end{remark}

\begin{remark}[Dotted circles are not RT colours]
\label{rem:dotted-circles-not-rt-colours}
A dotted circle in a one-ended dotted--framed Kirby record is notation for a
four-dimensional $1$-handle; it is not a coloured ribbon component and has
no intrinsic Reshetikhin--Turaev evaluation.  The RT evaluation above is
applied only after every dotted circle has been translated into an actual
categorical morphism.  That additional translation belongs to the
Crane--Yetter evaluator and is defined in
\cref{def:dotted-circle-translation}.  This is the dotted-circle notation of
\cite[Section~2.3.2 and Definition~2.49]{BarenzBarrett2018}.
\end{remark}

\begin{remark}[Graph evaluation versus a $3$-manifold invariant]
\label{rem:rt-graph-versus-manifold}
The evaluation in \cref{def:rt-graph-evaluation} only requires a ribbon
category.  It is not the surgery construction of an RT invariant of an
arbitrary closed $3$-manifold; the latter requires nondegeneracy and a
surgery normalization.
\end{remark}
